\chapter{Flujo de Ingresos Proyectados}

\section{Introducción}

La proyección del flujo de ingresos permite traducir la propuesta técnica del sistema en una estructura económico-comercial verificable. En una tesis aplicada, este capítulo cumple la función de responder quién compra la solución, qué paquete de servicio se le ofrece, cuánto se cobra, con qué periodicidad y bajo qué supuestos esos ingresos pueden sostenerse durante el horizonte de evaluación~\cite{europeancommission2014cba,unido2017manual}. Por ello, el análisis no puede limitarse a un flujo genérico por museo ni a una suma de hardware de bajo costo; debe reconocer que la solución combina diagnóstico institucional, despliegue BLE, estructuración curatorial, operación de IA, soporte y mejora continua.

El capítulo se apoya en tres premisas desarrolladas previamente. Primero, el mercado objetivo de la solución es institucional: los compradores directos son museos, unidades ejecutoras, municipalidades, universidades, fundaciones o redes culturales, aunque el usuario final sea el visitante. Segundo, el enfoque BYOD es viable porque OSIPTEL reportó que al cierre de 2024 el 94.8\% de hogares peruanos contaba con smartphone~\cite{osiptel2024smartphone}. Tercero, la oferta tecnológica comparada muestra que la solución propuesta no compite solamente como audioguía, sino como servicio de interpretación cultural digital con localización de sala, interacción por voz, recuperación aumentada y analítica de uso.

En consecuencia, el capítulo reformula el flujo de ingresos como un modelo B2B institucional, sustentado en compradores reales, paquetes comerciales diferenciados y escenarios financieros a diez años. El piloto del Museo Tumbas Reales de Sipán conserva valor como caso ancla, pero deja de ser la única unidad económica del análisis: se lo incorpora dentro de una cartera comercial más amplia, priorizada por demanda, capacidad institucional y necesidad de mediación digital.

\section{Mercado abordable y demanda institucional}

El mercado nacional no debe presentarse como una cifra homogénea de visitantes o como un conjunto indiferenciado de museos. La propia planificación sectorial reconoce limitaciones de medición: para museos, el registro sistemático se concentra principalmente en las instituciones administradas por el Ministerio de Cultura~\cite{cultura2020politica,cultura2025pei}. Esta advertencia metodológica obliga a separar datos observados de estimaciones del estudio.

Como aproximación al mercado total, la \textit{Guía de museos e instituciones museales del Perú} reúne información de 386 museos e instituciones museales del país~\cite{cultura2023guiamuseos}. Dentro de ese universo, el portal oficial del Ministerio de Cultura identifica 56 museos administrados por el sector, que constituyen el núcleo público más directamente articulable con políticas nacionales, unidades ejecutoras y compras institucionales~\cite{museoscultura}. Este submercado no agota la oportunidad comercial, pero sí ofrece una base prioritaria para validar pilotos y escalar estándares.

La demanda presencial refuerza la pertinencia del servicio. La serie oficial muestra una contracción severa durante la pandemia y una recuperación posterior: 353\,153 visitas en 2020, 1\,067\,571 en 2022, 1\,543\,872 en 2023 y alrededor de 1.79 millones en 2024 para museos administrados por el sector Cultura~\cite{cultura2020memoria,cultura2022memoria,cultura2023memoria,cultura2024reporte}. Para 2025, el informe institucional disponible reporta 1\,385\,895 visitantes nacionales a museos y evidencia concentración regional en Lambayeque, Lima, Cusco, Ayacucho y Áncash~\cite{cultura2025evaluacion}. Esta concentración geográfica orienta la prospección comercial hacia territorios donde existe flujo real de público y mayor probabilidad de justificar inversión.

\begin{table}[h!]
	\centering
	\footnotesize
	\setlength{\tabcolsep}{4pt}
	\resizebox{\textwidth}{!}{%
	\begin{tabular}{>{\raggedright\arraybackslash}p{5.7cm} c >{\raggedright\arraybackslash}p{6.1cm}}
		\hline
		\textbf{Indicador} & \textbf{Valor} & \textbf{Lectura para la solución} \\
		\hline
		Museos e instituciones museales identificados oficialmente & 386 & Mercado nacional amplio y heterogéneo, útil como techo de oportunidad \\
		Museos administrados por el Ministerio de Cultura & 56 & Núcleo público prioritario para pilotos, compras institucionales y escalamiento \\
		Hogares peruanos con smartphone al cierre de 2024 & 94.8\% & Condición habilitante para una arquitectura BYOD sin equipos de préstamo \\
		Visitas reportadas a museos del sector Cultura en 2024 & aprox. 1.79 millones & Masa crítica para justificar soluciones de mediación digital presencial \\
		Departamentos líderes en visitas nacionales 2025 & Lambayeque, Lima, Cusco, Ayacucho y Áncash & Prioridad geográfica para prospección y expansión \\
		\hline
	\end{tabular}
	}
	\caption[Síntesis del mercado abordable de la solución]{Síntesis del mercado abordable de la solución propuesta en el Perú. Fuente: elaboración propia con base en~\cite{cultura2023guiamuseos,museoscultura,osiptel2024smartphone,cultura2024reporte,cultura2025evaluacion}.}
	\label{tab:mercado-abordable-museiq}
\end{table}

Además de la demanda agregada, existe una necesidad sectorial explícita. El Ministerio de Cultura define el servicio de interpretación cultural como el conjunto de estrategias curatoriales, educativas y comunicativas que permiten generar experiencias significativas para el visitante~\cite{cultura2021interpretacion}. En la misma línea, la guía metodológica del MEF para proyectos de museos e instituciones museales públicas confirma que la mejora de servicios museales forma parte de una tipología legítima de inversión pública~\cite{mef2025guiamuseos}. Por tanto, la solución propuesta debe presentarse como una herramienta para interpretación cultural, accesibilidad y gestión de públicos, no como un accesorio tecnológico aislado.

\section{Compradores potenciales reales}

La prospección comercial debe partir de instituciones nominales y no de la categoría genérica ``museos peruanos''. Los compradores potenciales se priorizan por afluencia, valor patrimonial, complejidad narrativa, capacidad de articulación institucional y posibilidad de replicar la solución en una red o clúster territorial. Bajo ese criterio, la cartera comercial se amplía a 35 museos e instituciones museales nacionales. Esta ampliación mantiene el foco inicial en Lambayeque, Lima, Cusco, Ayacucho y Áncash, pero incorpora otros nodos del circuito público nacional para que el escenario base de diez años tenga compradores identificables y paquetes asignados.

\begin{table}[h!]
	\centering
	\scriptsize
	\setlength{\tabcolsep}{3pt}
	\resizebox{\textwidth}{!}{%
	\begin{tabular}{c >{\raggedright\arraybackslash}p{7.3cm} >{\raggedright\arraybackslash}p{2.4cm} c c}
		\hline
		\textbf{N.°} & \textbf{Institución objetivo} & \textbf{Región} & \textbf{Prioridad} & \textbf{Paquete sugerido} \\
		\hline
		1 & Museo Tumbas Reales de Sipán & Lambayeque & Muy alta & Avanzado \\
		2 & Museo de Sitio Pucllana & Lima & Muy alta & Avanzado \\
		3 & Museo Histórico Regional del Cusco & Cusco & Muy alta & Avanzado \\
		4 & Museo Nacional de Arqueología, Antropología e Historia del Perú & Lima & Muy alta & Avanzado \\
		5 & Museo Arqueológico y Etnográfico del Conjunto Monumental Belén & Cajamarca & Alta & Estándar \\
		6 & Museo de Sitio Pachacamac & Lima & Muy alta & Avanzado \\
		7 & Lugar de la Memoria, la Tolerancia y la Inclusión Social & Lima & Alta & Estándar \\
		8 & Museo de Sitio Wari & Ayacucho & Alta & Estándar \\
		9 & Museo de Sitio de Chan Chan & La Libertad & Alta & Avanzado \\
		10 & Museo Nacional Chavín & Áncash & Alta & Estándar \\
		11 & Museo Nacional Sicán & Lambayeque & Alta & Estándar \\
		12 & Museo de Sitio Túcume & Lambayeque & Alta & Estándar \\
		13 & Museo Arqueológico Nacional Brüning & Lambayeque & Alta & Estándar \\
		14 & Museo de Arte de Lima -- MALI & Lima & Muy alta & Avanzado \\
		15 & Museo Central -- MUCEN & Lima / red regional & Alta & Avanzado \\
		16 & Museo Pedro de Osma & Lima & Alta & Estándar \\
		17 & Museo de Artes y Tradiciones Populares Luis Repetto Málaga & Lima & Media alta & Básico / Estándar \\
		18 & Museo de Artes y Ciencias Ing. Eduardo de Habich & Lima & Media & Básico \\
		19 & Museo de Arqueología y Antropología de la UNMSM & Lima & Media alta & Estándar \\
		20 & Yuyanapaq. Para recordar & Lima & Media alta & Estándar \\
		21 & Museo Nacional del Perú -- MUNA & Lima & Muy alta & Avanzado \\
		22 & Museo Nacional de la Cultura Peruana & Lima & Alta & Estándar \\
		23 & Museo de Arte Italiano & Lima & Alta & Estándar \\
		24 & Museo de Sitio Arturo Jiménez Borja -- Puruchuco & Lima & Alta & Estándar \\
		25 & Museo de Sitio Huallamarca & Lima & Media alta & Básico / Estándar \\
		26 & Museo de Sitio Julio C. Tello de Paracas & Ica & Alta & Estándar \\
		27 & Museo Regional de Ica Adolfo Bermúdez Jenkins & Ica & Alta & Estándar \\
		28 & Museo de Sitio Manuel Chávez Ballón & Cusco & Muy alta & Avanzado \\
		29 & Museo de Sitio de Chinchero & Cusco & Alta & Estándar \\
		30 & Sala de Exhibición de Pikillaqta & Cusco & Media alta & Básico / Estándar \\
		31 & Sala de Exhibición de la Zona Arqueológica Monumental de Kotosh & Huánuco & Media alta & Estándar \\
		32 & Museo Histórico Regional Hipólito Unanue & Ayacucho & Alta & Estándar \\
		33 & Museo Histórico Regional de Tacna & Tacna & Media alta & Estándar \\
		34 & Museo Regional de Arqueología de Junín & Junín & Media & Básico / Estándar \\
		35 & Museo de Sitio de Narihualá & Piura & Media alta & Estándar \\
		\hline
	\end{tabular}%
	}
	\caption[Cartera ampliada de compradores potenciales]{Cartera ampliada de compradores potenciales para la solución propuesta. El paquete sugerido representa el alcance comercial inicial más probable según valor patrimonial, complejidad del recorrido, escala institucional y posibilidad de adopción. Fuente: elaboración propia con base en~\cite{museoscultura,cultura2023guiamuseos,infocultura2023visitantes,naylamp2023visitantes,cultura2025evaluacion}.}
	\label{tab:compradores-potenciales-museiq}
\end{table}

La Tabla~\ref{tab:compradores-potenciales-museiq} convierte el mercado en una cartera comercial concreta. En el corto plazo, Sipán, Pucllana, Cusco y el Museo Nacional de Arqueología funcionan como vitrinas de validación por su visibilidad y complejidad narrativa. En el mediano plazo, el clúster Lambayeque, Lima y los museos de sitio permiten replicar aprendizajes de instalación, corpus curatorial y soporte. En la segunda mitad del horizonte, la expansión hacia Ica, Huánuco, Tacna, Junín y Piura permite sostener el crecimiento sin asumir una penetración masiva del mercado nacional.

\section{Naturaleza económica y proceso comercial}

La propuesta debe entenderse como un sistema híbrido producto-servicio. Por un lado, incorpora elementos físicos: beacons BLE, nodos o dispositivos de borde cuando sean necesarios, señalética QR y accesorios de instalación. Por otro, concentra su valor en servicios profesionales y operación continua: diagnóstico de salas, diseño de cobertura, preparación de corpus curatorial, configuración de aplicación o PWA, integración RAG, capacitación, soporte, analítica y actualización de contenidos. Esta estructura es coherente con los sistemas producto-servicio, en los que el valor se entrega mediante la combinación de activos físicos y servicios sostenidos en el tiempo~\cite{tukker2004pss}.

La consecuencia económica es directa: la solución no debe monetizarse como venta de hardware ni como licencia genérica de software. El hardware BLE tiene bajo costo relativo, pero no produce valor sin calibración, contenidos y operación; el software requiere adaptación a cada museo para reconocer zonas, recuperar información curatorial pertinente y reportar métricas útiles. Por ello, la estrategia comercial combina pago inicial por implementación, contratos recurrentes y servicios adicionales.

El proceso comercial recomendado se organiza en ocho etapas: identificación del museo objetivo, evaluación preliminar de encaje, diagnóstico de salas y recorridos, propuesta técnica y económica por paquete, instalación de BLE y señalética, configuración de aplicación y corpus curatorial, pruebas con personal del museo, y operación recurrente con soporte, analítica y mejora continua. Esta secuencia alinea la propuesta con el servicio de interpretación cultural definido por el sector público y evita reducir el proyecto a una compra puntual de dispositivos~\cite{cultura2021interpretacion,mef2025guiamuseos}.

\begin{figure}[h!]
	\centering
	\includegraphics[width=0.82\textwidth]{estructura_ingresos_rubro.png}
	\caption[Estructura comercial de paquetes]{Estructura comercial de paquetes de la solución propuesta según precio de implementación y componentes recurrentes. Fuente: elaboración propia.}
	\label{fig:estructura-ingresos-museiq}
\end{figure}

\section{Paquetes de servicio}

Para volver defendible el flujo de ingresos, la oferta se organiza en tres paquetes: Básico, Estándar y Avanzado. Esta segmentación evita usar un precio único para instituciones muy distintas y permite adecuar alcance, complejidad y valor esperado. Los montos se expresan en soles constantes de 2026, sin IGV y sin financiamiento. Deben interpretarse como estimaciones razonadas para tesis, no como cotizaciones vinculantes.

\begin{table}[h!]
	\centering
	\scriptsize
	\setlength{\tabcolsep}{2.5pt}
	\resizebox{\textwidth}{!}{%
	\begin{tabular}{>{\raggedright\arraybackslash}p{2.0cm} >{\raggedright\arraybackslash}p{3.7cm} c c >{\raggedright\arraybackslash}p{6.1cm} c c c c}
		\hline
		\textbf{Paquete} & \textbf{Museo recomendado} & \textbf{Cobertura} & \textbf{BLE estimados} & \textbf{Funciones incluidas} & \textbf{Implementación} & \textbf{Mant. anual} & \textbf{IA y analítica anual} & \textbf{Costo directo impl.} \\
		\hline
		Básico & Museos municipales, universitarios, salas de memoria y museos pequeños & 3--5 salas & 8--12 & PWA o web app BYOD, QR, reconocimiento de sala, contenido textual e imagen, un idioma, analítica básica, soporte remoto y capacitación inicial & S/48\,000 & S/8\,500 & S/5\,500 & S/29\,000 \\
		Estándar & Museos medianos, museos de sitio y universidades con mayor flujo & 6--10 salas & 18--25 & Todo lo básico, más TTS, RAG curatorial acotado, dos idiomas, dashboard, reportes de uso, soporte preventivo y capacitación ampliada & S/96\,000 & S/18\,000 & S/11\,000 & S/58\,000 \\
		Avanzado & Museos nacionales, privados de alta afluencia, redes museales y espacios emblemáticos & 12--18 salas & 35--45 & Todo lo estándar, más STT/TTS, RAG integral, perfiles de recorrido, accesibilidad reforzada, hasta tres idiomas, SLA y mejora continua & S/185\,000 & S/34\,000 & S/18\,000 & S/111\,000 \\
		\hline
	\end{tabular}%
	}
	\caption[Paquetes comerciales propuestos]{Paquetes comerciales propuestos para la solución. Fuente: elaboración propia con base en el dimensionamiento técnico de la solución y la investigación económico-comercial del capítulo.}
	\label{tab:paquetes-museiq}
\end{table}

El paquete Básico reduce la barrera de entrada para museos pequeños o instituciones educativas. El paquete Estándar se considera la unidad comercial central para museos de sitio y museos medianos, porque combina reconocimiento de sala, voz, RAG acotado y analítica sin exigir la complejidad de una red completa. El paquete Avanzado corresponde a instituciones emblemáticas o redes museales con mayor demanda, narrativa extensa y necesidad de accesibilidad o multilingüismo reforzado.

La diferencia frente al cálculo microeconómico del piloto es importante. Un piloto académico puede instalar pocos sensores y usar una base de contenidos reducida, pero una implementación comercial debe financiar levantamiento, diseño de experiencia, limpieza de información, pruebas, capacitación, gestión contractual, soporte y contingencias. Por ello, el precio no se justifica por el costo unitario del beacon, sino por la solución completa de interpretación cultural digital.

\section{Fuentes de ingreso y servicios adicionales}

La estructura de ingresos de la solución propuesta se compone de cuatro rubros. El primero es la implementación inicial, cobrada por museo o sede. El segundo es el mantenimiento técnico anual, que sostiene continuidad operativa, revisiones, soporte, pequeñas actualizaciones y reposición menor. El tercero es el servicio anual de IA y analítica, que cubre operación de voz, recuperación aumentada, hosting, reportes e infraestructura digital. El cuarto corresponde a servicios adicionales, contratados según necesidad.

\begin{table}[h!]
	\centering
	\footnotesize
	\setlength{\tabcolsep}{4pt}
	\begin{tabular}{>{\raggedright\arraybackslash}p{4.2cm} >{\raggedright\arraybackslash}p{4.0cm} >{\raggedright\arraybackslash}p{6.1cm}}
		\hline
		\textbf{Fuente de ingreso} & \textbf{Periodicidad} & \textbf{Justificación económica} \\
		\hline
		Implementación inicial & Pago único por proyecto & Recupera diagnóstico, hardware, instalación, configuración, corpus inicial, pruebas y capacitación \\
		Mantenimiento técnico & Anual & Sostiene continuidad, soporte remoto, revisiones, ajustes menores y reposición básica \\
		IA y analítica & Anual, con posible sobreconsumo & Cubre operación de voz, RAG, hosting, tableros, monitoreo y reportes de uso \\
		Servicios adicionales & Bajo demanda & Permite ampliar idiomas, nuevas salas, exposiciones temporales, accesibilidad y formación especializada \\
		\hline
	\end{tabular}
	\caption[Fuentes de ingreso del modelo comercial]{Fuentes de ingreso del modelo comercial de la solución propuesta. Fuente: elaboración propia.}
	\label{tab:fuentes-ingreso-museiq}
\end{table}

\begin{table}[h!]
	\centering
	\footnotesize
	\setlength{\tabcolsep}{4pt}
	\begin{tabular}{>{\raggedright\arraybackslash}p{7.2cm} >{\raggedright\arraybackslash}p{5.2cm}}
		\hline
		\textbf{Servicio adicional} & \textbf{Precio referencial} \\
		\hline
		Nueva sala, zona o exposición temporal & S/4\,000--S/12\,000 según alcance \\
		Idioma adicional o módulo narrativo multilingüe & S/5\,000--S/18\,000 \\
		Accesibilidad ampliada, audiodescripción o rutas inclusivas & S/6\,000--S/20\,000 \\
		Capacitación adicional para personal del museo & S/2\,500--S/8\,000 \\
		Integración con CMS, catálogo, ticketing o portal institucional & Según alcance técnico \\
		Campaña educativa o turística asociada al lanzamiento & S/5\,000--S/15\,000 \\
		\hline
	\end{tabular}
	\caption[Servicios adicionales plausibles]{Servicios adicionales plausibles para la solución propuesta. Fuente: elaboración propia.}
	\label{tab:servicios-adicionales-museiq}
\end{table}

El componente de IA merece una separación explícita porque su costo depende del uso. En el escenario comercial, la suscripción anual debe incluir una bolsa de uso coherente con el paquete contratado y una regla de sobreconsumo cuando la interacción por voz supere el umbral acordado. Esta política evita que una mayor adopción del visitante deteriore el margen del proveedor. La lógica ya fue observada en el piloto de Sipán: si una proporción relevante de visitantes usa STT, el costo operativo deja de ser marginal y debe monetizarse como servicio recurrente.

\section{Supuestos de proyección a diez años}

La proyección se construye sobre un horizonte de diez años en soles constantes de 2026. No se incluye IGV, deuda, depreciación financiera ni subsidios, porque el objetivo del capítulo es formular el flujo económico-comercial de ingresos y costos operativos directos. Los nuevos contratos se reconocen en promedio a mitad de año; por ello, los ingresos de mantenimiento e IA de clientes nuevos se computan al 50\% durante el año de captación y al 100\% desde el año siguiente.

El escenario base asume una captación gradual de 35 museos en diez años. Esta cifra equivale a menos del 10\% del universo nacional de 386 instituciones identificado por la guía oficial y no exige cubrir todo el sistema público desde el inicio. La mezcla total de la cartera nominal considera 5 contratos Básicos, 20 Estándar y 10 Avanzados. En los casos marcados como Básico/Estándar, la proyección adopta el paquete Básico como entrada prudente y deja la ampliación posterior como servicio adicional o mejora contractual.

\begin{table}[h!]
	\centering
	\footnotesize
	\setlength{\tabcolsep}{4pt}
	\resizebox{\textwidth}{!}{%
	\begin{tabular}{>{\raggedright\arraybackslash}p{8.7cm} c}
		\hline
		\textbf{Supuesto} & \textbf{Valor adoptado} \\
		\hline
		Horizonte de proyección & 10 años \\
		Moneda & Soles constantes de 2026 \\
		Impuestos y financiamiento & No incluidos \\
		Reconocimiento recurrente de clientes nuevos & 50\% en el año de captación \\
		Paquetes considerados & Básico, Estándar y Avanzado \\
		Museos captados en escenario base & 35 al año 10 \\
		Opex fijo inicial & S/130\,000 en el año 1 \\
		Crecimiento anual de opex fijo & 5\% aproximado \\
		Servicios adicionales & Entre 8\% y 10\% de implementación, según mezcla anual \\
		\hline
	\end{tabular}
	}
	\caption[Supuestos de proyección comercial]{Supuestos de proyección comercial de la solución propuesta. Fuente: elaboración propia.}
	\label{tab:supuestos-proyeccion-museiq}
\end{table}

\section{Escenario base detallado}

El escenario base plantea una expansión progresiva desde casos ancla hacia una cartera de instituciones activas. Para responder a la observación metodológica sobre el carácter genérico del flujo, la proyección no se presenta únicamente como cantidad de museos nuevos, sino como un cronograma nominal de ventas. Cada año identifica a qué institución se vendería la solución propuesta y qué paquete se considera como primera contratación. Esta asignación no representa una carta de intención ni una venta comprometida; es una ruta comercial razonada para sustentar el flujo de ingresos a diez años.

\begin{table}[h!]
	\centering
	\scriptsize
	\setlength{\tabcolsep}{3pt}
	\resizebox{\textwidth}{!}{%
	\begin{tabular}{c >{\raggedright\arraybackslash}p{8.4cm} c >{\raggedright\arraybackslash}p{4.6cm}}
		\hline
		\textbf{Año} & \textbf{Instituciones captadas y paquete de entrada} & \textbf{B/E/A} & \textbf{Lógica comercial del año} \\
		\hline
		1 & Museo Tumbas Reales de Sipán (Avanzado); Museo de Sitio Pucllana (Avanzado) & 0/0/2 & Casos ancla de alta visibilidad en Lambayeque y Lima para demostrar valor patrimonial, turístico y operativo. \\
		2 & Museo Histórico Regional del Cusco (Avanzado); Museo Nacional de Arqueología, Antropología e Historia del Perú (Avanzado) & 0/0/2 & Consolidación con museos nacionales o regionales emblemáticos y relatos curatoriales extensos. \\
		3 & Museo Arqueológico y Etnográfico del Conjunto Monumental Belén (Estándar); Museo de Sitio Pachacamac (Avanzado); Lugar de la Memoria, la Tolerancia y la Inclusión Social (Estándar) & 0/2/1 & Entrada a museos de sitio y memoria con recorridos interpretativos complejos. \\
		4 & Museo de Sitio Wari (Estándar); Museo de Sitio de Chan Chan (Avanzado); Museo Nacional Chavín (Estándar) & 0/2/1 & Expansión hacia nodos arqueológicos de alto valor patrimonial fuera de Lima. \\
		5 & Museo Nacional Sicán (Estándar); Museo de Sitio Túcume (Estándar); Museo Arqueológico Nacional Brüning (Estándar) & 0/3/0 & Profundización del clúster Lambayeque para aprovechar cercanía territorial y gestión sectorial afín. \\
		6 & Museo de Arte de Lima -- MALI (Avanzado); Museo Central -- MUCEN (Avanzado); Museo Pedro de Osma (Estándar); Museo de Artes y Tradiciones Populares Luis Repetto Málaga (Básico) & 1/1/2 & Ingreso a instituciones privadas, financieras y universitarias con mayor flexibilidad de adopción. \\
		7 & Museo de Artes y Ciencias Ing. Eduardo de Habich (Básico); Museo de Arqueología y Antropología de la UNMSM (Estándar); Yuyanapaq. Para recordar (Estándar); Museo Nacional del Perú -- MUNA (Avanzado) & 1/2/1 & Combinación de museos universitarios, memoria y un museo nacional de alto potencial. \\
		8 & Museo Nacional de la Cultura Peruana (Estándar); Museo de Arte Italiano (Estándar); Museo de Sitio Arturo Jiménez Borja -- Puruchuco (Estándar); Museo de Sitio Huallamarca (Básico) & 1/3/0 & Cobertura de museos limeños medianos y de sitio con paquetes escalables. \\
		9 & Museo de Sitio Julio C. Tello de Paracas (Estándar); Museo Regional de Ica Adolfo Bermúdez Jenkins (Estándar); Museo de Sitio Manuel Chávez Ballón (Avanzado); Museo de Sitio de Chinchero (Estándar); Sala de Exhibición de Pikillaqta (Básico) & 1/3/1 & Expansión sur hacia Ica y Cusco, combinando museos de sitio y salas con distinta complejidad. \\
		10 & Sala de Exhibición de la Zona Arqueológica Monumental de Kotosh (Estándar); Museo Histórico Regional Hipólito Unanue (Estándar); Museo Histórico Regional de Tacna (Estándar); Museo Regional de Arqueología de Junín (Básico); Museo de Sitio de Narihualá (Estándar) & 1/4/0 & Cierre del horizonte con museos regionales que consolidan presencia nacional sin exigir cobertura total del sistema. \\
		\hline
	\end{tabular}%
	}
	\caption[Cronograma nominal de captación del escenario base]{Cronograma nominal de captación del escenario base. B/E/A indica la cantidad anual de contratos Básico, Estándar y Avanzado. Fuente: elaboración propia con base en la cartera ampliada de compradores potenciales.}
	\label{tab:cronograma-nominal-museiq}
\end{table}

La Tabla~\ref{tab:cronograma-nominal-museiq} asigna nombre propio al flujo proyectado. La secuencia inicia con instituciones de alta visibilidad, continúa con clústeres territoriales y luego incorpora museos regionales, universitarios o de sitio. Así, el escenario base deja de depender de una variable abstracta de ``museos nuevos'' y se convierte en una cartera de ventas defendible.

\begin{table}[h!]
	\centering
	\scriptsize
	\setlength{\tabcolsep}{2.2pt}
	\resizebox{\textwidth}{!}{%
	\begin{tabular}{c c c c c c c c c c c c}
		\hline
		\textbf{Año} & \textbf{Nuevos} & \textbf{Activos} & \textbf{B/E/A} & \textbf{Ing. impl.} & \textbf{Ing. mant.} & \textbf{Ing. IA} & \textbf{Ing. adic.} & \textbf{Costo impl. y adic.} & \textbf{Costo mant. e IA} & \textbf{Opex fijo} & \textbf{Flujo neto} \\
		\hline
		1  & 2 & 2  & 0/0/2 & 370\,000 & 34\,000  & 18\,000  & 29\,600 & 236\,800 & 24\,000  & 130\,000 & 60\,800 \\
		2  & 2 & 4  & 0/0/2 & 370\,000 & 102\,000 & 54\,000  & 29\,600 & 236\,800 & 72\,000  & 136\,500 & 110\,300 \\
		3  & 3 & 7  & 0/2/1 & 377\,000 & 171\,000 & 92\,000  & 30\,160 & 242\,080 & 122\,000 & 143\,325 & 162\,755 \\
		4  & 3 & 10 & 0/2/1 & 377\,000 & 241\,000 & 132\,000 & 30\,160 & 242\,080 & 174\,000 & 150\,491 & 213\,589 \\
		5  & 3 & 13 & 0/3/0 & 288\,000 & 303\,000 & 168\,500 & 23\,040 & 185\,520 & 221\,000 & 158\,016 & 218\,004 \\
		6  & 4 & 17 & 1/1/2 & 514\,000 & 377\,250 & 211\,250 & 41\,120 & 329\,560 & 276\,500 & 165\,917 & 371\,643 \\
		7  & 4 & 21 & 1/2/1 & 425\,000 & 463\,750 & 260\,250 & 34\,000 & 273\,000 & 340\,500 & 174\,213 & 395\,287 \\
		8  & 4 & 25 & 1/3/0 & 336\,000 & 534\,250 & 302\,250 & 26\,880 & 216\,440 & 394\,500 & 182\,924 & 405\,516 \\
		9  & 5 & 30 & 1/3/1 & 521\,000 & 613\,750 & 349\,750 & 41\,680 & 334\,840 & 455\,500 & 192\,070 & 543\,770 \\
		10 & 5 & 35 & 1/4/0 & 432\,000 & 702\,250 & 402\,750 & 34\,560 & 278\,280 & 523\,500 & 201\,674 & 568\,106 \\
		\hline
	\end{tabular}%
	}
	\caption[Escenario base de flujo económico-comercial]{Escenario base de flujo económico-comercial de la solución propuesta a diez años, ajustado al cronograma nominal de la Tabla~\ref{tab:cronograma-nominal-museiq}. B/E/A indica la mezcla anual de paquetes Básico, Estándar y Avanzado. Fuente: elaboración propia.}
	\label{tab:flujo-base-museiq}
\end{table}

\begin{figure}[h!]
	\centering
	\includegraphics[width=0.82\textwidth]{flujo_ingresos_10_anios.png}
	\caption[Ingresos, costos y flujo neto del escenario base]{Ingresos, costos y flujo neto anual del escenario base. Fuente: elaboración propia.}
	\label{fig:flujo-ingresos-base-museiq}
\end{figure}

El resultado principal es que el modelo no depende indefinidamente de la implementación inicial. En los primeros años, la venta de paquetes impulsa el ingreso; desde la mitad del horizonte, la base instalada produce un volumen creciente de mantenimiento, IA y analítica. En el escenario base nominal, los ingresos acumulados ascienden a S/9\,863\,800 y los costos acumulados a S/6\,814\,030, generando un flujo neto acumulado aproximado de S/3\,049\,770 al final del año 10.

\section{Comparación de escenarios}

Para reconocer la incertidumbre institucional, se modelan tres escenarios. El conservador asume una entrada lenta, más dependiente de contratación pública, con 20 museos captados en diez años. El base incorpora 35 museos nominales y combina clústeres públicos con instituciones privadas, universitarias y municipales. El optimista llega a 49 museos, suponiendo mayor tracción por referencias, redes institucionales y contratos múltiples.

\begin{table}[h!]
	\centering
	\footnotesize
	\setlength{\tabcolsep}{3pt}
	\resizebox{\textwidth}{!}{%
	\begin{tabular}{>{\raggedright\arraybackslash}p{3.0cm} c c c c c c}
		\hline
		\textbf{Escenario} & \textbf{Museos captados} & \textbf{Activos año 10} & \textbf{Ingreso acumulado} & \textbf{Costo acumulado} & \textbf{Flujo neto acumulado} & \textbf{Equilibrio acumulado} \\
		\hline
		Conservador & 20 & 20 & S/4\,341\,000 & S/3\,957\,126 & S/383\,874 & Año 9 \\
		Base nominal ajustado & 35 & 35 & S/9\,863\,800 & S/6\,814\,030 & S/3\,049\,770 & Año 1 \\
		Optimista & 49 & 49 & S/11\,811\,000 & S/7\,883\,126 & S/3\,927\,874 & Año 1 \\
		\hline
	\end{tabular}%
	}
	\caption[Comparación de escenarios de flujo económico-comercial]{Comparación de escenarios de flujo económico-comercial de la solución propuesta a diez años. Fuente: elaboración propia.}
	\label{tab:escenarios-flujo-museiq}
\end{table}

\begin{figure}[h!]
	\centering
	\includegraphics[width=0.82\textwidth]{instituciones_activas_10_anios.png}
	\caption[Museos activos según escenario comercial]{Museos activos acumulados según escenario comercial. Fuente: elaboración propia.}
	\label{fig:instituciones-activas-museiq}
\end{figure}

\begin{figure}[h!]
	\centering
	\includegraphics[width=0.82\textwidth]{escenarios_ingresos_comparados.png}
	\caption[Ingresos anuales según escenario comercial]{Ingresos anuales según escenario comercial de la solución propuesta. Fuente: elaboración propia.}
	\label{fig:escenarios-ingresos-museiq}
\end{figure}

La comparación muestra que la sostenibilidad depende menos del margen de una instalación aislada y más de la capacidad de acumular instituciones activas. El escenario conservador sigue siendo positivo, pero recién alcanza equilibrio acumulado en la parte final del horizonte. El escenario base nominal ajustado es más apropiado para la tesis porque no exige penetración masiva del mercado nacional: captar 35 museos en diez años implica atender una fracción moderada del universo identificado oficialmente, priorizando nodos con mayor afluencia y capacidad institucional.

\section{Análisis crítico}

El principal avance del modelo es que transforma el flujo de ingresos en una cartera económico-comercial concreta. La solución ya no se presenta como un proyecto genérico de bajo costo, sino como un servicio institucional que debe venderse a compradores reales y sostenerse con paquetes diferenciados. Esta lógica mejora la defensa académica porque conecta mercado objetivo, oferta tecnológica, costos de implementación, operación recurrente y escenarios de adopción.

La primera limitación es estadística. La serie de visitas no cubre de manera homogénea todos los museos públicos, privados, municipales y universitarios; además, algunos reportes diferencian visitantes nacionales de visitantes totales. Por ello, las cifras de mercado deben leerse como evidencia de oportunidad y priorización, no como una base cerrada de demanda nacional. La segunda limitación es comercial: los precios de paquetes son estimaciones razonadas y requieren validación con entrevistas, cartas de interés o cotizaciones comparables. La tercera es operativa: la rentabilidad depende de controlar el costo variable de IA, especialmente STT, TTS y procesamiento conversacional, mediante bolsas de uso y sobrecargos transparentes.

También existe riesgo de adopción. En museos públicos, los ciclos de contratación pueden ser lentos y la continuidad presupuestal no siempre está garantizada. En museos privados, la decisión puede ser más rápida, pero la exigencia de retorno en experiencia, reputación o ticketing será mayor. En ambos casos, la capacitación del personal y la actualización del corpus curatorial son condiciones críticas: sin gobernanza de contenidos, la solución puede degradarse aunque la infraestructura técnica funcione.

Pese a estos límites, el modelo resulta más realista que una proyección basada únicamente en el piloto de Sipán o en la venta de beacons. Su fortaleza reside en separar compradores, paquetes, ingresos únicos, recurrencias, costos operativos y escenarios de adopción. Esa separación prepara una evaluación financiera posterior más defendible y permite discutir la sostenibilidad de la solución como servicio tecnológico cultural.

\section{Conclusión}

El flujo de ingresos del proyecto debe estructurarse como un modelo B2B institucional orientado a servicios de interpretación cultural digital. El mercado abordable existe y cuenta con evidencia suficiente: 386 museos e instituciones museales identificados oficialmente, 56 museos administrados por el Ministerio de Cultura, recuperación de visitas presenciales y alta disponibilidad de smartphones para sostener una arquitectura BYOD. Sin embargo, la oportunidad solo se vuelve comercialmente defendible cuando se priorizan compradores reales y se ofrecen paquetes adecuados a la escala de cada institución.

La propuesta se organiza en tres paquetes: Básico, Estándar y Avanzado, con precios de implementación de S/48\,000, S/96\,000 y S/185\,000, respectivamente, más mantenimiento, IA y analítica anual. Bajo el escenario base nominal ajustado, la captación gradual de 35 museos en diez años genera ingresos acumulados de S/9\,863\,800 y un flujo neto acumulado aproximado de S/3\,049\,770. En consecuencia, el modelo deja de ser un flujo genérico y se convierte en una estructura económico-comercial sustentada en instituciones compradoras, servicios claros y una proyección financiera defendible.
